% Residencial II - apartamento com varanda, duas suites e acessos coerentes
\section{Projeto Residencial II — apartamento com varanda}

\newcommand{\AptoFinalBase}{%
\draw[arqExt] (0,0) rectangle (12,9);\draw[arqExt] (12,4) rectangle (14,9);
\ArqParedeInt{4.2}{0}{4.2}{9}\ArqParedeInt{5.2}{2}{5.2}{6.2}\ArqParedeInt{0}{4.4}{4.2}{4.4}
\ArqParedeInt{2.6}{0}{2.6}{2.1}\ArqParedeInt{2.6}{2.1}{4.2}{2.1}
\ArqParedeInt{2.6}{6.8}{4.2}{6.8}\ArqParedeInt{2.6}{6.8}{2.6}{9}
\ArqParedeInt{5.2}{3.8}{12}{3.8}\ArqParedeInt{7.2}{0}{7.2}{3.8}\ArqParedeInt{10}{0}{10}{3.8}
\ArqParedeInt{5.2}{2.2}{6.6}{2.2}\ArqParedeInt{6.6}{2.2}{6.6}{3.8}
% entrada e acessos
\ArqPortaHUp{5.7}{0}{1.0}\node[arqSub] at (6.2,0.45){ENTRADA / HALL COMUM};
\ArqVaoH{5.65}{7.05}{3.8} % foyer -> sala
\ArqVaoH{7.65}{9.45}{3.8} % cozinha -> sala
\ArqPortaVLeft{5.2}{2.45}{0.85} % foyer -> hall intimo
\ArqPortaVLeft{4.2}{2.55}{0.85}\ArqPortaVLeft{4.2}{4.90}{0.85} % suites
% banhos compactos: portas de correr evitam conflito com loucas e box
\ArqPortaCorrerH{2.95}{3.75}{2.1}\ArqPortaCorrerH{2.95}{3.75}{6.8}
\ArqPortaHUp{5.45}{2.2}{0.75} % lavabo abre para o foyer
\ArqPortaVRight{10}{1.10}{0.80} % cozinha -> servico
\ArqPortaCorrerV{5.0}{8.2}{12} % sala -> varanda
% janelas
\ArqJanV{0.55}{3.4}{0}\ArqJanV{5.25}{8.35}{0}\ArqJanH{7.65}{9.4}{0}\ArqJanH{10.35}{11.6}{0}\ArqJanH{0.6}{2.3}{9}
% suite 2
\ArqCama{0.45}{0.50}{1.95}{2.55}\ArqArmario{0.4}{3.45}{2.8}{0.5}\ArqChuveiro{2.78}{0.15}{1.15}{0.78}\ArqVaso{3.55}{1.35}\ArqPia{2.72}{1.15}{0.55}{0.4}\node[arqLabel] at (1.55,3.9){SUÍTE 2};\node[arqSub] at (3.45,1.92){B2};
% suite master
\ArqCama{0.45}{4.95}{2.05}{2.70}\ArqArmario{0.45}{8.05}{1.85}{0.5}\ArqChuveiro{2.78}{8.05}{1.15}{0.78}\ArqVaso{3.55}{7.35}\ArqPia{2.72}{7}{0.55}{0.4}\node[arqLabel] at (1.5,8.55){SUÍTE MASTER};\node[arqSub] at (3.45,8.72){B1};
\node[arqSub,rotate=90] at (4.70,4.25){HALL ÍNTIMO};
% foyer e lavabo
\node[arqSub] at (6.2,1.45){FOYER};\ArqVaso{5.75}{2.85}\ArqPia{5.35}{3.25}{0.55}{0.38}\node[arqSub] at (5.9,3.6){LAVABO};
% cozinha e servico
\ArqBancada{7.45}{0.25}{2.15}{0.55}\ArqBancada{9.35}{0.8}{0.55}{2.5}\ArqFogao{8.55}{0.18}\ArqGeladeira{7.35}{2.7}\node[arqLabel] at (8.55,3.35){COZINHA};
\ArqLavadora{10.35}{0.35}\ArqBancada{11.1}{0.3}{0.55}{2.65}\node[arqLabel] at (11.05,3.35){SERVIÇO};
% estar e varanda
\ArqSofa{5.8}{4.55}{3.2}{1.05}\ArqBancada{5.55}{7.75}{0.38}{1.0}\ArqMesa{8.65}{6.35}{2.15}{1.0}\node[arqLabel] at (8.65,5.95){ESTAR / JANTAR};
\ArqSofa{12.45}{5.05}{1.1}{2.2}\ArqMesa{12.55}{7.65}{0.85}{0.7}\node[arqLabel,rotate=90] at (13.55,6.5){VARANDA};
\CotaH{0}{12}{9.18}{12,00 m}\CotaV{0}{9}{14.18}{9,00 m}
}

\begin{frame}{Residencial II — acessos refinados}
\centering\begin{tikzpicture}[scale=0.53,every node/.style={font=\tiny}]\AptoFinalBase\end{tikzpicture}
\vspace{0.4mm}\scriptsize Entrada pelo hall comum; lavabo voltado ao foyer; cozinha integrada à sala; serviço acessado pela cozinha; B1/B2 por dentro das suítes, com portas de correr; varanda pela sala.
\end{frame}

\begin{frame}{Residencial II — iluminação A1/A2}
\centering\begin{tikzpicture}[scale=0.53,every node/.style={font=\tiny}]
\AptoFinalBase\SimQD{6.85}{3.30}{QD-A}
\foreach \x/\y in {1.45/2.8,3.35/1,1.5/6,3.35/8,4.7/4.3,6.2/1.2,5.9/2.9,8.5/1.9,11/1.9,8/5.1,10.1/6.5,13/6.6}{\SimLuz{\x}{\y}}
\foreach \x/\y in {3.92/2.82,3.12/2.32,3.92/5.12,3.12/6.58,5.42/0.42,5.42/2.55,6.82/3.52,9.82/1.28,5.42/4.18,11.72/5.18}{\SimInterruptor{\x}{\y}{S}}
\draw[corAccent!80!black,dashed,thick] (6.85,3.30)--(8.0,5.1)--(10.1,6.5)--(13,6.6);
\draw[corAccent!80!black,dashed,thick] (6.85,3.30)--(8.5,1.9)--(11,1.9);
\draw[corAccent!55!black,dashed,thick] (6.85,3.30)--(6.2,1.2)--(5.9,2.9)--(4.7,4.3)--(1.5,6)--(3.35,8);
\draw[corAccent!55!black,dashed,thick] (4.7,4.3)--(1.45,2.8)--(3.35,1);
\node[font=\tiny\bfseries,fill=white,inner sep=1pt,corAccent!80!black] at (9.5,4.8){A1};
\node[font=\tiny\bfseries,fill=white,inner sep=1pt,corAccent!55!black] at (3.2,4.3){A2};
\end{tikzpicture}
\vspace{0.3mm}\scriptsize Previsão: A1 = 0,94 kVA e A2 = 0,94 kVA. Os comandos estão junto aos acessos; a potência mínima total é 1,88 kVA, calculada pelas áreas e pelos acréscimos completos de 4 m$^2$.
\end{frame}

\begin{frame}{Residencial II — 27 TUG e sete cargas específicas}
\centering\begin{tikzpicture}[scale=0.53,every node/.style={font=\tiny}]
\AptoFinalBase\SimQD{6.85}{3.30}{QD-A}
% A3: suites, hall, foyer, sala e varanda — 16 pontos
\foreach \x/\y in {0.25/0.55,0.25/4.15,2.35/0.25,3.95/4.15,0.25/4.85,0.25/8.75,2.35/8.75,3.95/4.65,4.35/4.30,5.35/0.55,5.45/4.05,7.40/4.05,9.50/4.05,11.70/4.10,11.70/8.70,13.70/8.70}{\SimTUG{\x}{\y}}
% A4 cozinha, A5 servico, A6 banheiros
\foreach \x/\y in {7.45/0.25,8.55/0.25,9.75/1.50,7.35/3.55}{\SimTUGM{\x}{\y}}
\foreach \x/\y in {10.15/0.25,11.75/0.25,10.15/3.55,11.75/3.55}{\SimTUGM{\x}{\y}}
\foreach \x/\y in {3.95/1.65,3.95/7.15,5.45/3.50}{\SimTUG{\x}{\y}}
% A7--A13
\SimTUE{3.4}{0.4}{A7}\SimTUE{3.4}{8.55}{A8}\SimTUE{8.9}{0.35}{A9}\SimTUE{10.7}{0.45}{A10}
\SimTUE{0.3}{3.1}{A11}\SimTUE{0.3}{7.95}{A12}\SimTUE{11.55}{7.8}{A13}
\end{tikzpicture}
\vspace{0.3mm}\tiny A3 áreas secas/varanda (16) · A4 cozinha (4) · A5 serviço (4) · A6 banheiros/lavabo (3) · A7--A8 chuveiros · A9 forno · A10 lavadora · A11--A13 ar-condicionado.
\end{frame}

\begin{frame}[shrink=8]{Residencial II — memória dos mínimos por ambiente}
\scriptsize
\begin{tabularx}{\textwidth}{>{\bfseries}p{2.35cm}c c c c >{\centering\arraybackslash}X}
\toprule
Ambiente & Área (m$^2$) & Perím. (m) & Luz (VA) & TUG & TUG (VA)\\
\midrule
Suíte 2 / B2 & 15,1 / 3,4 & 17,2 / próprio & 220 / 100 & 4 / 1 & 400 / 600\\
Suíte master / B1 & 15,8 / 3,5 & 17,6 / próprio & 220 / 100 & 4 / 1 & 400 / 600\\
Hall íntimo & 4,2 & critério por área & 100 & 1 & 100\\
Foyer & 5,4 & critério por área & 100 & 1 & 100\\
Lavabo & 2,2 & critério próprio & 100 & 1 & 600\\
Sala/jantar & 35,4 & 24,0 & 520 & 5 & 500\\
Cozinha & 10,6 & 13,2 & 160 & 4 & 1.900\\
Serviço & 7,6 & 11,6 & 100 & 4 & 1.900\\
Varanda & 10,0 & ao menos uma & 160 & 1 & 100\\
\midrule
\textbf{Total} & -- & -- & \textbf{1.880} & \textbf{27} & \textbf{7.200}\\
\bottomrule
\end{tabularx}
\end{frame}

\begin{frame}[shrink=6]{Residencial II — circuitos A1 a A6}
\scriptsize
\begin{tabularx}{\textwidth}{c p{3.5cm} c c c c c X}
\toprule
C & Uso & $S$ & $I_b$ & Cu & DJ & Fase & Solução adotada\\
\midrule
A1 & Luz social/cozinha/serviço/varanda & 0,94 kVA & 4,3 A & 1,5 mm$^2$ & 10 A & A & RCBO 30 mA\\
A2 & Luz íntima/foyer/lavabo & 0,94 kVA & 4,3 A & 1,5 mm$^2$ & 10 A & B & RCBO 30 mA\\
A3 & TUG áreas secas e varanda & 1,60 kVA & 7,3 A & 2,5 mm$^2$ & 16 A & C & RCBO 30 mA\\
A4 & TUG cozinha & 1,90 kVA & 8,6 A & 2,5 mm$^2$ & 16 A & B & RCBO 30 mA\\
A5 & TUG serviço & 1,90 kVA & 8,6 A & 2,5 mm$^2$ & 16 A & C & RCBO 30 mA\\
A6 & TUG banheiros/lavabo & 1,80 kVA & 8,2 A & 2,5 mm$^2$ & 16 A & C & RCBO 30 mA\\
\bottomrule
\end{tabularx}
\end{frame}

\begin{frame}[shrink=6]{Residencial II — circuitos A7 a A13}
\scriptsize
\begin{tabularx}{\textwidth}{c p{3.1cm} c c c c c X}
\toprule
C & Equipamento & $P$ & $I_b$ & Cu & DJ & Fase & Proteção/observação\\
\midrule
A7 & Chuveiro B2 & 4,50 kW & 20,5 A & 4 mm$^2$ & 25 A & A & dedicado; RCBO 30 mA\\
A8 & Chuveiro B1 & 4,50 kW & 20,5 A & 4 mm$^2$ & 25 A & B & dedicado; RCBO 30 mA\\
A9 & Forno elétrico & 3,50 kW & 15,9 A & 4 mm$^2$ & 20 A & C & dedicado; RCBO 30 mA\\
A10 & Lavadora & 1,20 kW & 5,5 A & 2,5 mm$^2$ & 16 A & B & dedicado; RCBO 30 mA\\
A11--A13 & Ar-condicionado & 3$\times$1,00 kW & 4,5 A cada & 2,5 mm$^2$ & 16 A & A/A/A & dedicados; fabricante\\
\bottomrule
\end{tabularx}
\begin{caixaAt}
Seleções iniciais: recalcular $I_z$, agrupamento, queda, curto, capacidade de interrupção e requisitos do fabricante. A proteção diferencial adotada deve ser compatível com cargas eletrônicas/inverter.
\end{caixaAt}
\end{frame}

\begin{frame}{Residencial II — balanceamento e interface EMUC}
\small\begin{columns}[T,onlytextwidth]
\column{0.49\textwidth}
\begin{caixaProjeto}[title=Unidade revisada]
\begin{align*}
S_A&=\SI{8.44}{kVA}\;(\SI{38.4}{A})\\
S_B&=\SI{8.54}{kVA}\;(\SI{38.8}{A})\\
S_C&=\SI{8.80}{kVA}\;(\SI{40.0}{A})
\end{align*}
Total equivalente: \SI{25.78}{kW}. Proteção de 63 A é a premissa inicial da unidade. Para esta ilustração, os VA de previsão foram tomados numericamente como W; a declaração real usa potências ativas nominais. Verificar também cenários de demanda, pois o balanceamento instalado concentra os três aparelhos de ar-condicionado na fase A.
\end{caixaProjeto}
\column{0.49\textwidth}
\begin{caixaNorma}[title=Edifício multifamiliar]
Centro de medição, demanda do conjunto, prumadas, queda, proteção geral e solução de conexão são definidos pela \textbf{NT.00004.EQTL Rev.08} e seu Anexo I vigente. Não multiplicar a carga desta unidade pelo número de apartamentos sem aplicar o método do empreendimento.
\end{caixaNorma}
\end{columns}
\end{frame}

\begin{frame}{Residencial II — unifilar da unidade}
\centering\begin{tikzpicture}[scale=0.74,every node/.style={font=\tiny}]
\tikzset{b/.style={draw=corPrim,fill=corDef!8,rounded corners,minimum width=2cm,minimum height=0.56cm,align=center}}
\node[b](cm)at(0,4.4){Centro de medição\\EMUC};\node[b](pr)at(0,3.5){Prumada 3F+N+PE};\node[b](qd)at(0,2.55){QD-A\\DG 3P 63 A + DPS};
\node[b](l)at(-3.6,0.9){A1--A6\\luz/TUG + RCBO};\node[b](ch)at(-1.2,0.9){A7--A8\\chuveiros};\node[b](co)at(1.2,0.9){A9--A10\\forno/lavadora};\node[b](ac)at(3.6,0.9){A11--A13\\climatização};
\node[b,draw=corNorma,fill=corNorma!7](bep)at(4.8,2.9){BEP do edifício\\PE/equipot.};
\draw[-{Stealth},thick,corPrim](cm)--(pr)--(qd);\foreach \x in {l,ch,co,ac}{\draw[-{Stealth},thick](qd)--(\x);}\draw[corNorma,thick](pr)--(bep);\draw[corNorma,thick](qd)--(bep);\draw[corNorma,thick](bep.south)--(4.8,1.55)node[ground]{};
\end{tikzpicture}\end{frame}
